\documentclass[a4paper]{article}
% Español (México) con XeLaTeX: fontspec para fuentes Unicode, polyglossia
% para la división silábica y los nombres del documento en la variante mexicana.
\usepackage{fontspec}
\usepackage{polyglossia}
\setdefaultlanguage[variant=mexican]{spanish}
% Los nombres chinos van como «pinyin (original)»: xeCJK da a los caracteres
% Han su propia fuente; sin ella XeLaTeX los omite con un aviso.
% FandolSong no cubre algunos caracteres de nombres (U+73FA); WenQuanYi Zen Hei sí.
\usepackage[AutoFallBack=true]{xeCJK}
\setCJKmainfont{FandolSong-Regular.otf}
\setCJKfallbackfamilyfont{\CJKrmdefault}{WenQuanYi Zen Hei}
% polyglossia no traduce los nombres de listings.
\AtBeginDocument{\providecommand{\lstlistingname}{}\renewcommand{\lstlistingname}{Listado}%
  \providecommand{\lstlistlistingname}{}\renewcommand{\lstlistlistingname}{Índice de listados}}
\usepackage{amsmath, amssymb}
\usepackage{graphicx}
\usepackage[margin=2.5cm]{geometry}
\usepackage[most]{tcolorbox}
\usepackage{etoolbox}
\usepackage{listings}
\usepackage{hyperref}
\usepackage{booktabs}
\usepackage{float}
\newtcolorbox{knowledgebox}[1]{enhanced, colback=blue!5!white, colframe=blue!75!black, colbacktitle=blue!75!black, coltitle=white, fonttitle=\bfseries, title={#1}, attach boxed title to top left={yshift=-2mm, xshift=2mm}, boxrule=1pt, sharp corners}
\newtcolorbox{importantbox}[1]{enhanced, colback=yellow!10!white, colframe=yellow!80!black, colbacktitle=yellow!80!black, coltitle=black, fonttitle=\bfseries, title={#1}, sharp corners}
\newtcolorbox{warningbox}[1]{enhanced, colback=red!5!white, colframe=red!75!black, colbacktitle=red!75!black, coltitle=white, fonttitle=\bfseries, title={#1}, sharp corners}
\lstset{language=Python, basicstyle=\ttfamily\small, keywordstyle=\color{blue}, stringstyle=\color{red!60!black}, commentstyle=\color{green!60!black}, breaklines=true, frame=single, numbers=left, numberstyle=\tiny\color{gray}, captionpos=b, extendedchars=true}
\newcommand{\notetitle}{Detalles de cómputo del Agentic Loop de veRL: asincronía y gestión de estado}
\newcommand{\noteauthors}{Elaboradas a partir de materiales públicos del curso}
\newcommand{\notedate}{2025}
\newcommand{\videochannel}{Wudaokou Nash (五道口纳什)}
\newcommand{\videopublishdate}{2025}
\newcommand{\videoduration}{aproximadamente 20 minutos}
\newcommand{\videourl}{}
\newcommand{\videocoverpath}{cover.jpg}
\newcommand{\repourl}{https://github.com/hqhq1025/ai-course-notes}

% Footer with repo info
\usepackage{fancyhdr}
\pagestyle{fancy}
\fancyhf{}
\fancyfoot[C]{\thepage}
\fancyfoot[R]{\footnotesize\color{gray}\href{\repourl}{AI Course Notes}}
\renewcommand{\headrulewidth}{0pt}
\renewcommand{\footrulewidth}{0pt}

\begin{document}
\begin{titlepage}\centering
{\Large Notas del curso\par}\vspace{1.2cm}{\huge\bfseries \notetitle\par}\vspace{0.8cm}{\large \noteauthors\par}\vspace{0.3cm}{\large \notedate\par}\vspace{1.2cm}
\ifdefempty{\videocoverpath}{}{\includegraphics[width=0.82\textwidth,height=0.45\textheight,keepaspectratio]{\videocoverpath}\par}
\vfill
\begin{tcolorbox}[width=0.9\textwidth, colback=black!2!white, colframe=black!60, sharp corners]
\textbf{Autor o canal del video}: \videochannel\par\textbf{Fecha de publicación}: \videopublishdate\par\textbf{Duración del video}: \videoduration\par
\ifdefempty{\videourl}{}{\textbf{Enlace del video}: \href{\videourl}{\nolinkurl{\videourl}}\par}\par
\textbf{Página del proyecto}: \href{\repourl}{\nolinkurl{\repourl}}\par
\end{tcolorbox}
\end{titlepage}
\tableofcontents\newpage

\section{Introducción}
En esta sección se responden preguntas clave sobre la implementación del Agentic Loop: ¿por qué debe ser asíncrono? ¿cómo se gestionan los distintos estados?

\section{¿Por qué asincronía?}
\begin{importantbox}{La necesidad de la asincronía}
En el Agentic Loop, distintos sample tienen un número diferente de pasos de Multi-turn. Si se espera de forma síncrona al sample más lento, la eficiencia es extremadamente baja. La ejecución asíncrona permite que los sample ya completados regresen primero, maximizando la utilización de la GPU.
\end{importantbox}

\section{Gestión de estados}
En el Agentic Loop, cada sample puede estar en diferentes estados:
\begin{itemize}
  \item \textbf{Generating}: está generando texto
  \item \textbf{Tool Calling}: está ejecutando una herramienta
  \item \textbf{Waiting}: espera a que la herramienta regrese
  \item \textbf{Done}: ha completado todas las rondas
\end{itemize}

La transición de estados se controla mediante el archivo de configuración: si se configura una tool, se llama a la tool; si se configura una interaction, se simula la retroalimentación del usuario.

\subsection{Resumen de la sección}
La ejecución asíncrona y una gestión de estados fina son fundamentales para el funcionamiento eficiente del Agentic Loop.

\newpage
\section{Recordatorio de ingeniería: entre más estados, más se necesitan límites explícitos}
La capacidad de ejecución asíncrona del Agentic Loop depende, en esencia, de una máquina de estados. Si la definición de los estados no es clara o las condiciones de transición son ambiguas, el sistema cae fácilmente en bloqueos, ejecuciones repetidas o una recuperación incorrecta de recursos a lo largo de las llamadas a herramientas en múltiples rondas.

\begin{warningbox}{Riesgos comunes en la gestión de estados}
\begin{itemize}
  \item El sample tarda demasiado en llegar a Done, lo que provoca fugas en la cola
  \item Después de que la herramienta responde, no se regresa correctamente a Generating
  \item Hay demasiados parámetros de configuración, pero la lógica de transición no tiene un punto de entrada unificado
\end{itemize}
\end{warningbox}

\subsection{Resumen de la sección}
Lo que más teme un sistema asíncrono es que «los estados no estén claros». Definir con precisión los límites de los estados y las condiciones de transición es más importante que perseguir sin criterio una mayor concurrencia.

\newpage
\section{Síntesis y ampliación}
\begin{enumerate}
  \item La ejecución asíncrona resuelve el problema de eficiencia que surge porque distintos samples tienen un número diferente de pasos
  \item La gestión de estado asegura que cada sample ejecute la operación correcta en la etapa correcta
  \item El diseño impulsado por configuración ofrece flexibilidad
\end{enumerate}
\end{document}
